\documentclass[11pt]{article}
\usepackage[a4paper,margin=25mm,headheight=15pt]{geometry}
\usepackage[T1]{fontenc}
\usepackage{lmodern,amsmath,amssymb,amsthm,mathtools,booktabs,longtable,array,microtype}
\usepackage{xcolor,enumitem,fancyhdr}
\usepackage[hidelinks,breaklinks]{hyperref}
\usepackage{xurl}
\hypersetup{pdftitle={Compact proof certificates for ZFC consistency Round 13},pdfauthor={Prepared for Vladimir Reshetnikov},pdfsubject={Explicit Diophantine equations and reproducible ordinary mathematical reductions}}
\setlength{\parindent}{0pt}
\setlength{\parskip}{5pt plus 1pt}
\setlength{\emergencystretch}{3em}
\setlist{nosep,leftmargin=1.7em}
\pagestyle{fancy}\fancyhf{}
\fancyhead[L]{\small Compact proof certificates}
\fancyhead[R]{\small Round 13}
\fancyfoot[C]{\thepage}
\renewcommand{\headrulewidth}{0pt}
\newtheorem{theorem}{Theorem}[section]
\newtheorem{lemma}[theorem]{Lemma}
\newtheorem{proposition}[theorem]{Proposition}
\newtheorem{corollary}[theorem]{Corollary}
\DeclareMathOperator{\Con}{Con}
\begin{document}
\begin{center}
{\LARGE\bfseries Compact proof certificates for ZFC consistency\par}
\vspace{5pt}
{\large Round 13 research report and reproducible equations\par}
\vspace{10pt}
Prepared for Vladimir Reshetnikov\\
24 October 2026
\end{center}
\vspace{4pt}
\begin{abstract}
We give a smaller explicit direct proof-certificate construction for a
Diophantine formulation of the consistency of ZFC. The main equation is a
quartic in $1859$ nonnegative integer variables; a distinct signed-quotient
companion is a quartic in $3179$ unrestricted integer variables. Their
coefficients have absolute value below $2^{26}$ and are supplied in full.
The main variable reduction is from the independently recounted pinned
$3443$-variable baseline, not an obsolete predecessor. The construction removes
local beta remainder slacks, replaces copied bound circuits by fixed
squaring chains, and makes canonical Pell witnesses point forward before
eliminating their definitional coordinates. Its final synthesis composes
these reductions with incoming eight-column projection and sign-orbit
normalization, with those imported ideas credited explicitly.

A separate inherited native-code route now yields a twelve-integer
equation of exact degree $2048$ with $233$ arithmetic operations. Its
new quintic-carrier synthesis uses a supplied full-residue interpolation
theorem. Additional complete equations have variable/degree pairs
$13/1072$, $14/378$, $15/172$, $15/170$ and $16/98$. Finally, four ten-integer norm circuits become shorter, while a
leading-form argument proves degree $1080888$ optimal only within the
specified final-coordinate translation family. None of these results is a
global minimum or a proof that ZFC is consistent. The report distinguishes
ordinary mathematical derivations, inherited calibration, exact finite
polynomial checks, and the end-to-end formalization still to be completed.
\end{abstract}

\section{Results and the exact logical target}
\label{sec:target}
Fix the recursive ZFC axiom presentation and finite proof predicate used
by the included calibration sources. Write
\[
 \Con(\mathrm{ZFC})\;:=\;
 \neg\exists p\in\mathbb N\;\operatorname{Ref}_{\mathrm{ZFC}}(p).
\]
Here $\mathbb N=\{0,1,2,\ldots\}$ throughout. For every proposed polynomial
one must specify its coefficient tuple and its domain; a count over
$\mathbb N$ is not interchangeable with a count over $\mathbb Z$.
The intended statement is a theorem of ZFC of the form
\[
 \mathrm{ZFC}\vdash
 \bigl(\neg\exists x\in D^n\;P(x)=0\bigr)
 \ \longleftrightarrow\ \Con(\mathrm{ZFC}),
 \qquad D\in\{\mathbb N,\mathbb Z\}.
\]
Equivalently, the construction must internally identify existence of a
zero with a finite refutation. It must not assume consistency in proving
that identification. Nor does the identification prove either side of
the displayed equivalence. The finite identities and witness
transformations given here are ordinary arithmetic arguments formalizable
in ZFC. Their composition uses the named inherited refutation calibration.
The package is not a new kernel-checked proof of the complete display.

The comparison snapshot is
\href{https://github.com/VladimirReshetnikov/ProveIt/tree/7e533c33819ee1f703b69f2778a767cdc069940d}{ProveIt commit 7e533c33819ee1f703b69f2778a767cdc069940d},
authored at 03:32:10 UTC on 24 October 2026. It is a fixed snapshot, not
an assertion that the repository has stopped changing.

\begin{table}[htbp]
\centering\small
\begin{tabular}{llrrrr}
\toprule
Construction & Domain & Variables & Degree & Terms & Height\\
\midrule
Pinned direct baseline & $\mathbb N$ &3443&4&15754&127302193\\
New direct variable focus & $\mathbb N$ &1859&4&12275&63924066\\
New direct sparse option & $\mathbb N$ &2560&4&11074&63923744\\
New direct integer companion & $\mathbb Z$ &3179&4&23157&63923766\\
\bottomrule
\end{tabular}
\caption{Complete collected quartics. Height is the maximum absolute
integer coefficient, after collecting equal monomials.}
\end{table}

For the other route, literal complete graph counts are as follows.
A circuit operation is one live binary addition, subtraction or
multiplication, including a square or scalar multiplication. Literal
integer leaves and reuse are free. An abbreviation or circuit gate is
not an additional unknown.
\begin{center}
\begin{tabular}{lrrr}
\toprule
Family & Integer variables & Exact degree & Operations\\
\midrule
Quintic five-root carrier &12&2048&233\\
Five-root coprime carrier &13&1072&192\\
Four-root coprime carrier &14&378&145\\
Three-root direct carrier &15&172&123\\
Three-root centered carrier &15&170&134\\
Two-root centered carrier &16&98&138\\
Ten-integer polar norm &10&1245324&467\\
Ten-integer linear translation &10&1179788&470\\
Ten-integer cubic translation &10&1114252&473\\
Ten-integer six-term translation &10&1080888&507\\
\bottomrule
\end{tabular}
\end{center}
These rows cannot be ranked by a single unstated meaning of ``simplest.''
The direct route has small coefficient height and degree but many
variables. The native route has very few variables and large fixed
numerals. The centered fifteen-variable point trades eleven extra
operations for two degrees relative to its direct sibling. We claim
only the stated, audited tradeoffs.

\subsection{What is new and what is inherited}
The direct construction uses the numerical C222 proof calculus, its
ZF-refutation calibration, and constructible-choice passage from ZF to
ZFC refutations. It is independent of the K215/Jones/native coefficient
route. The low-variable families retain the latter route's coherent
coefficient tuple and its ordinary calibration. Keeping the two chains
separate is essential: a short evaluator identity cannot replace an
unbounded compiler or Pell theorem.

The local work in this round consists of explicit source transformations,
new carrier combinations, a credited projection and sign-normalization
synthesis, polynomial identities, new exact circuits,
and a restricted-family lower bound. The baseline comparison and all
coefficient payloads are fixed before counts are made. Source files are
retained with notices as inert mathematical provenance; no upstream
executable is run by the release checks. The checkers supply finite
implementation evidence, not a proof of consistency or a replacement
for universal soundness and completeness.

\subsection{How an explicit equation is specified here}
For each direct equation the file \texttt{equation.json} gives an ordered
list of variables and quadratic rows with literal integer coefficients.
A row is encoded as a finite list $(c,\mu)$, where $\mu$ is a list of
zero-based variable indices and denotes $c\prod_{i\in\mu}x_i$.
The empty list denotes the constant monomial. The output is the sum of
the squares of all rows. The accompanying
\texttt{expanded\_polynomial.json} gives its collected integer polynomial,
not just a hash or a recipe for missing coefficients.

For a small-variable norm, the JSON specifies a finite acyclic graph with
only literal or fixed-integer leaves and binary integer arithmetic. Every
fixed-integer leaf has its actual byte payload in the archive, together
with signedness, byte order and hashes. Hence these are ordinary fixed
polynomials, not coefficient-parameter families. The factored displays
below define the same equations more legibly than a huge monomial listing.


\subsection{Comparison update from later incoming packages}
After the fixed repository snapshot was reviewed, additional manuscripts
and equation packets arrived. They change the comparison, but not the
literal correctness of the equations in this report. They are not silently
treated as files that were already present at the pinned commit.

A separate literal audit re-expanded the complete sums of squares in
those direct packets and checked every supplied collected coefficient.
Among them are unrestricted-integer quartics with the following metrics:
\begin{center}
\begin{tabular}{rrrr}
\toprule
Variables & Quadratic rows & Collected terms & Height\\
\midrule
3776&2015&20194&47959\\
3671&1942&26271&3835369234\\
\bottomrule
\end{tabular}
\end{center}
The initial $4087$-integer checkpoint was therefore not a new smallest
integer count among the incoming packets. The final targeted synthesis in
Section~\ref{sec:projected-direct} instead gives $3179$ unrestricted
integers, fewer than both rows. It also improves the second row's
monomial count and height, but not the first row's small coefficients or
monomial count. The $1859$-natural-variable construction is smaller in
dimension than the later $2473$-natural quartic, whose literal audit gives
$2089$ rows, $20580$ terms and height $127808289$. Projection and sign
normalization from these incoming packages are credited prerequisites of
that final synthesis, not claimed as independently new ideas here.
These are domain-specific comparisons, not a total ordering.

A later thirteen-integer graph has $190$ live binary operations and a
claimed exact degree $1072$, compared with $192$ operations here. The
operation count was independently verified; its cancellation-sensitive
exact degree remains supported by its accompanying ordinary proof rather
than that finite recount alone. Consequently the thirteen-variable row
here is not presented as a current operation record.
Two later twelve-integer graphs have $248$ and $223$ operations and
reported degrees $4832$ and $16000$, respectively. Their syntactic degree
upper bounds and operation counts were checked; exact attainment and
the unbounded consistency calibration are separate authored claims.
They provide further variable--degree tradeoffs rather than invalidating
the smaller-degree thirteen- to sixteen-variable equations here.

The comparison addendum, audit receipts and payload identities are retained under
\texttt{late\_review/}. The original comparison archives and their separate checkers are not
part of the compact offline suite; these receipts are comparison evidence,
not a promise to reproduce omitted packages. The literal audit is not an
independent reproof of all inherited calibrations or a claim that those
later equations were derived in this round. We keep the principal
$3443\to1859$ improvement explicitly tied to the stated fixed baseline.

\section{Compressed direct proof-certificate quartics}
\label{sec:round13-direct}

The final selected direct construction gives a literal quartic in 1,859
nonnegative variables and a literal quartic in 3,179 unrestricted integer
variables.  The natural equation has 1,581 quadratic residuals, 12,275
monomials and coefficient height 63,924,066.  The integer equation has
2,074 residuals, 23,157 monomials and height 63,923,766.  Both heights are
below $2^{26}$.  For its integer polynomial $Q_8$, the theorem is
\[
 \operatorname{Con}(\mathrm{ZFC})
 \quad\Longleftrightarrow\quad
 \neg\exists z\in\mathbb Z^{3179}\;Q_8(z)=0.
\]
This is a no-zero equivalence, not a proof that the polynomial has no zero,
and not a new kernel-checked end-to-end translation.

We first establish the arithmetic and an eleven-column intermediate
source, then project to eight columns and normalize auxiliary signs in
Section~\ref{sec:projected-direct}.  The reconstructible intermediate
checkpoint has 2,170 natural variables or 4,087 integer variables; these
are not the final selected counts.  A retained sparse alternative with
2,560 natural variables and 11,074 monomials is also supplied.  Its smaller
monomial count is a genuine tradeoff against the selected 1,859-variable
natural polynomial.  The already improved baseline before this round has
3,443 natural variables, 2,838 residuals, 15,754 monomials and coefficient
height 127,302,193.  Late incoming integer comparisons and their advantages
are stated at the end of this section.

The construction follows the direct numerical C222 proof calculus,
independently of universal-polynomial or small-machine routes.  Its
calibration and indexed-Pell theorem are explicit inherited dependencies.
All new transformations below are existence-preserving mathematical
arguments, and the final files contain literal integer coefficients.

\subsection{The exact direct source and the carrier argument}

Write $I(x,y)=2(x+y)(x+y+1)+4y$, and define
\[
 A(x,y)=I(x,y)+2,\quad M(x,y)=I(x,y)+1,\quad
 E(x,y)=I(x,y)+3,\quad \neg x=I(x,1).
\]
These are numerical constructors, not arithmetic assertions about sets.
The four schema coordinates $a,b,c,d$ and twelve padded Horn branches are
exactly the pinned numerical C222 source.  The target is code $1$ and the
free truth code is $0$.  The explicit axiom constructors, including the
bundled axiom and its disequality guard, occur literally in
\texttt{build\_compact.py}; the inherited sparse source and calibration
modules are retained in \texttt{sources/}.

The dependency is the C222-to-ZF refutation theorem followed by the
constructible-choice ZF-to-ZFC refutation theorem, not execution of a
small universal machine.  In particular, no assumption that every natural
number decodes a well-formed formula is silently added.

For $N,m\ge1$, let $M=(m+1)^N$ and bound each rule label $p_i,q_i,v_i$
by $m$.  The accumulator source is
\[
 S_0=M+1,\quad S_N=1,\qquad
 S_{i+1}(1+v_iM)=S_i(1+p_iM)(1+q_iM),\quad i<N.
 \tag{1}
\]
Every row also requires the pinned Horn rule $(p_i,q_i)\to v_i$.
Telescoping gives an equality at $M$ between
\[
 F(X)=\prod_{i<N}(1+v_iX),\qquad
 G(X)=(1+X)\prod_{i<N}(1+p_iX)(1+q_iX).
\]
If $\deg G\ge N+1$, nonnegative coefficients give
$G(M)\ge M^{N+1}$, whereas
\[
 F(M)\le(1+mM)^N<((m+1)M)^N=M^{N+1}.
\]
Thus both products have at most $N$ nonconstant factors.  Their coefficient
sums at $1$ are at most $(m+1)^N=M$, and every individual coefficient is
strictly below $M$: a nonconstant product has its positive constant term
and another positive coefficient, while the constant product is $1<M$.
A nonzero polynomial $F-G$ cannot vanish at $M$, since its least nonzero
coefficient would be a nonzero multiple of $M$ with absolute value below
$M$.  Hence $F=G$.

The distinct factors $1+jX$, $j>0$, now recover the exact nonnegative
Horn flow balance for target $1$.  A finite single-conclusion flow implies
derivability: if the target were not derivable, sum the flow coordinates
outside the deductive closure.  Each valid rule contributes at most zero,
but the target contributes one.  Conversely, expand a finite proof tree
from the open obligation $1$ and let $S_i$ be the product of the weights
of all open obligations.  This gives (1).  Truth padding has weight one.
The proof uses single conclusions and nonnegative multiplicities; it does
not justify arbitrary cyclic proofs or signed flows.

\subsection{Eleven coordinates instead of thirteen}

Use the beta sequence function
\[
 \beta(x,y,i)=x\bmod\bigl(1+(i+1)y\bigr).
\]
For every finite natural sequence, its scale $y$ may be chosen strictly
larger than every sequence entry and divisible by $N!$.  The moduli for
$i=0,\ldots,N$ are then pairwise coprime: a common divisor is coprime to
$y$ and divides a positive index difference at most $N$, which divides
$y$.  The Chinese remainder theorem supplies $x$.

Parameterize
\[
 N=N_0+1,\quad m=m_0+1,\quad M=M_0+1=(m+1)^N,
 \qquad y=M_0+2+y_0,
\]
with all displayed underlying parameters natural.  The two endpoint rows
are simply
\[
 x=M_0+2+(y+1)u_0,
 \qquad x=1+\bigl(1+(N+1)y\bigr)u_N. \tag{2}
\]
They are canonical remainder statements because $M_0+2\le y<y+1$ and
$1<1+(N+1)y$.  The parameterization alone permits $S_0=y$; completeness
chooses $y>S_0$, and the first bounded row also enforces this strict
inequality.  No incorrect strict endpoint inequality is assumed.

At row $i$, retain only
\[
 w=(p,q,v,a,b,c,d,S,S',u,u')\in\mathbb N^{11}.
\]
Besides the Horn condition and the accumulator row, require
\[
 x=S+\bigl(1+(i+1)y\bigr)u,
 \qquad x=S'+\bigl(1+(i+2)y\bigr)u'. \tag{3}
\]
The CRT compiler bounds $p,q,v$ by $m$, bounds $S,S'$ by $y-1$, and
bounds $a,b,c,d,u,u'$ by one independent $U$.  Therefore (3) already
forces the canonical beta digits.  The two local remainder slack
coordinates and their two binomial range certificates disappear.

For completeness first choose the proof, $N,m,M$, and its finite
accumulator sequence, then choose $y$ strictly above every state, then
choose $x$ by CRT, and finally choose $U$ above the finitely many schema
and quotient coordinates.  None of these choices depends on the later
bounded-universal CRT encoding.  This order excludes circular bounds.

\subsection{Factoring the common Horn conditions}

Let $f_1,\ldots,f_8$ be the eight unconditional axiom codes, let $b_{\rm bundle}$ denote
the bundled axiom code, and let
\[
 g=((c-a)(c-b))^2-(d+1)^2
\]
be its guard residual.  Here $b$ in the guard remains the schema coordinate.
Put
\[
 J=v\prod_{r=1}^{8}(v-f_r),\qquad
 K=(v-b_{\rm bundle})^2+g^2.
\]
The following fixed nonnegative polynomial has exactly the pinned Horn
zero set:
\[
 H_{\rm Horn}=(q-I(p,v))^2
 \left[q^2+(v-A(a,p))^2(p^2+J^2)(p^2+K)\right]. \tag{4}
\]
Indeed its first factor is modus ponens.  Otherwise the bracket vanishes
exactly when $q=0$ and at least one of these alternatives holds:
$v=A(a,p)$; or $p=0$ and $J=0$; or $p=0$ and $K=0$.
These are generalization, truth or one of the eight unconditional axioms,
and the guarded bundled axiom, respectively.  Every factor used in this
argument is nonnegative.  This is an equivalence of zero sets, not an
asserted polynomial identity with the old twelve-factor source.

Let $P$ be (4) plus the squares of the accumulator residual and the two
residuals in (3).  The signed arithmetic circuit is split as $P=E-F$,
where $E,F$ have nonnegative coefficients, using
\[
 (e,f)+(g,h)=(e+g,f+h),\qquad
 (e,f)(g,h)=(eg+fh,eh+fg),
\]
and $(e-f)^2=(e^2+f^2)-2ef$.  Thus $E=F$ holds exactly when every
required local condition holds.

\subsection{One small circuit for a very large majorant}

It is unnecessary to evaluate a second full copy of the source circuit at
all coordinate upper bounds.  Put
\[
 R=4+N+x+y+m+M_0+U.
\]
Every input at which the upper source is evaluated is at most $R$: the
index is $N$, the label bounds are $m$, the six other private bounds are
$U$, and $y$ is an upper bound for each state (slightly larger than its
actual bound $y-1$).  The external inputs are also included in $R$.

A structural exponent certificate bounds every nonnegative circuit node
by a power of $R$.  An input receives exponent one.  A constant $c$
receives the least $e$ with $c\le4^e$.  At addition take
$\max(e_1,e_2)+1$, since $R\ge4\ge2$; at multiplication take $e_1+e_2$.
Induction on the circuit proves the bound at every node.  For the actual
factored source the root exponents are 6,893,154 and 6,893,148.
Rounding upwards to $2^{23}=8,388,608$ gives
\[
 H=R^{2^{23}}\ge N,\qquad H\ge E(N,z,u),F(N,z,u). \tag{5}
\]
Here the upper state entries in $u$ may be taken to be $y$, as just noted.
The literal compiler computes $H$ by exactly 23 successive squarings.
There is no variable exponent or power-representation macro at this step.
The larger numerical value of the bound has essentially no circuit cost.
The exported exponent certificate records every abstract circuit node
and every one of these 23 squarings.

\subsection{One bounded-quantifier certificate with direct cancellation}

Set
\[
 T=H!,\quad D_i=1+(i+1)T,\quad B=\prod_{i<N}D_i.
\]
For actual coordinate bounds
$u=(m,m,m,U,U,U,U,y-1,y-1,U,U)$, the bounded row family is equivalent to
existence of positive encoded columns $A_j$ and a natural encoded index $C$
with
\[
 B\mid {A_j\choose u_j+1},\qquad
 B\mid1+(C+1)T,\qquad E(C,z,A)\equiv F(C,z,A)\pmod B. \tag{6}
\]
In particular, the two state binomial lower arguments are exactly $y$;
no subtraction is needed in the finite arithmetic circuit.

For soundness fix $i<N$ and take a prime $p\mid D_i$.  It is coprime to
$T=H!$, hence $p>H$.  The identity
\[
 (u_j+1)!{A_j\choose u_j+1}=\prod_{r=0}^{u_j}(A_j-r)
\]
implies that $A_j\equiv w_j\pmod p$ for some $0\le w_j\le u_j$.
This remains true when $A_j\le u_j$, as one factor is zero.  The index
condition gives $C\equiv i\pmod p$.  The evaluation congruence then holds
for $E(i,z,w),F(i,z,w)$; both lie in $[0,H]$, so they are equal.

For completeness the $D_i$ are pairwise coprime because $H\ge N$.
The radius bound supplies more than needed for soundness: $H\ge R$ and
$R\ge u_j+1$ for every coordinate.  Consequently
$L_j=(u_j+1)!$ divides $H!=T$, and $\gcd(B,L_j)=1$.
Choose ordinary CRT representatives
$A_j\equiv w_j^{(i)}\pmod{D_i}$ and $C\equiv i\pmod{D_i}$.
The falling-factorial numerator then vanishes modulo every $D_i$ and
hence modulo $B$.  Cancelling $L_j$, which is coprime to $B$, gives
$B\mid {A_j\choose u_j+1}$.  Adding a sufficiently large multiple of
$B$ makes $A_j>0$, preserving the falling-product congruence and hence,
by the same cancellation, the binomial divisibility.  The valid rows
and polynomial congruence give the evaluation condition in (6).

This final radius construction therefore needs neither prime-supported
CRT lifts nor their proof-only products.  Those are necessary in the more
general current-baseline theorem without a factorial-size guard, but the
very cheap radius bound restores the simpler direct cancellation proof.
There is no circularity: all coordinate bounds and $H$ are chosen before
$T,B$ and the encoded columns.

\subsection{All arithmetic operators are eliminated}

The exact arithmetic identities are
\[
 Q_0=(1+(N+1)T)^N+2,\qquad
 B=\left(T^N N!{Q_0+N\choose N}\right)\bmod(TQ_0-1), \tag{7}
\]
\[
 {n\choose k}=
 \left\lfloor\frac{(2^n+2)^n}{(2^n+1)^k}\right\rfloor\bmod(2^n+1), \tag{8}
\]
\[
 r=(n+2)^{n+2},\quad C_0={r+n\choose n},\qquad
 n!=\left\lfloor\frac{(r+1)^n+C_0-1}{C_0}\right\rfloor. \tag{9}
\]
For (7), reduction of the numerator modulo $TQ_0-1$ gives
$\prod_{j=1}^N(1+jT)$, and the chosen $Q_0$ makes this product strictly
smaller than the modulus.  Division uses $a=dq+r$ and $d=r+s+1$.
The circuit for (8) combines its two divisions using one high-digit
coordinate.  Formula (9) uses the equivalent equation
$(r+1)^n+C_0=C_0 n!+r_0+1$, with $0\le r_0<C_0$.
The included factorial and bounded-forall representation sources state
these identities with their boundary cases.

There are exactly 49 positive-power graph calls, 14 binomial calls, and
two factorial calls.  Positivity is established before each call:
$N,m\ge1$; $H\ge4$; encoded columns are parameterized as $A_j=\alpha_j+1$;
range lower arguments are $m+1$, $U+1$, or $y\ge2$; all auxiliary
factorial/binomial upper arguments and power bases are positive.
Thus deleting zero-base and zero-exponent cases is licensed for this
construction, not for the unrestricted exponentiation relation.

\subsection{The parametric positive-power graph}

For positive $n,k$, let $m$ be the proposed output.  There are twelve
natural certificate coordinates
\[
 a_2,y_k,g,d,X,U,q_4,q_B,q_S,q_R,q_P,Z.
\]
Define, rather than allocate additional unconstrained coordinates,
\begin{align*}
 W&=n+k,& A&=a_2+2,& Y&=k+y_k,& T_0&=m+g+1,\\
 V&=Y^2(d+1),& B_0&=1+4Y(q_4+1),&
 S&=X+Uq_S,& R_0&=k+4Yq_R.
\end{align*}
Writing $\mathcal N(a,x,y)$ for
$x^2+y^2-a^2y^2=1$, require precisely seven polynomial conditions:
\begin{gather*}
 \mathcal N(A,X,Y),\quad \mathcal N(A,U,V),\quad
 \mathcal N(B_0,S,R_0),\\
 B_0=A+Uq_B,\qquad
 X+Yn=YA+m+T_0q_P,\qquad 2An=T_0+n^2+1,\\
 \mathcal N(W+1,A,WZ). \tag{10}
\end{gather*}
These are circuit expressions; flattening their products gives quadratic
residuals.  The explicitly assigned $B_0$ is strictly greater than one,
including when $q_4=0$.  There is no missing bound on this Pell parameter.

\paragraph{Soundness.}
The definitions give $A>1$, $Y\ge k>0$, $V>0$, $Y^2\mid V$, and
$B_0>1$.  The three forward congruence equations imply the three modular
conditions of the pinned indexed-Pell criterion.  Thus $(X,Y)$ is the
$k$th Pell solution for parameter $A$.  The last norm gives
\[
 A^2=1+W^3(W+2)Z^2.
\]
As $A\ge2$ and $W\ge1$, necessarily $Z\ge1$ and $A>W\ge n$.
Therefore the rearranged power congruence uses the same ordinary
subtraction as the pinned power theorem.  Also $m<T_0$, and the deleted
common-bound slacks can be restored as $W-n=k$ and $W-k=n$.
The positive-power criterion yields $m=n^k$.

\paragraph{Completeness and the direction of every quotient.}
Let $(X_a(j),Y_a(j))$ denote the natural Pell sequence.  At $W=n+k$ choose
\[
 A=X_{W+1}(W),\quad WZ=Y_{W+1}(W),\quad
 T_0=2An-n^2-1,\quad X=X_A(k),\ Y=Y_A(k).
\]
The standard congruence $Y_{W+1}(W)\equiv W\pmod W$ supplies $Z$.
The usual fixed-bound power construction gives $n^k<T_0$ and the power
congruence.  For its direction set
$F_j=X_A(j)-(A-n)Y_A(j)$.  Then
$F_0=1,F_1=n,F_{j+1}=2AF_j-F_{j-1}$.  Inductively
$F_j\ge F_{j-1}\ge0$ and
$F_{j+1}\ge(2A-1)F_j\ge nF_j$, so $F_k\ge n^k$.
Thus $q_P=(F_k-n^k)/T_0$ is natural.

In the explicit completeness proof of the indexed-Pell criterion set
$h=2kY$, $U=X_A(h)$, $V=Y_A(h)$, and choose $B_0$ by CRT with
$B_0\equiv A\pmod U$ and $B_0\equiv1\pmod{4Y}$.
That proof has $A<U$, so its natural CRT representative satisfies
$B_0\ge A$.  Take $S=X_{B_0}(k)$ and $R_0=Y_{B_0}(k)$, at the same
index $k$.  The Pell sequences are monotone in the parameter, hence
$S\ge X$; also $R_0\ge k$.  Consequently all $q_B,q_S,q_R$ are natural.
Because $B_0>1$ and $Y>0$, its quotient in $B_0=1+4Yq$ satisfies $q\ge1$,
so write $q=q_4+1$.  The remaining slacks follow from the displayed
inequalities and divisibility.  This proves completeness of (10), not
merely its one-way soundness.  The separate orientation supplement
records the monotonicity proof and exact pinned source lines.

\subsection{Quartic flattening, sparse alternative, and integer domain}

Every nonnegative arithmetic addition or multiplication receives a
triangular defining row.  Nonnegative affine gate values may be
substituted without raising residual degree; the implementation substitutes
only expressions of support at most four, a size tradeoff rather than a
semantic restriction.  The complete substitution map is supplied.
The final polynomial is the sum of squares of all remaining quadratic
rows.  Over either $\mathbb N$ or $\mathbb Z$, a sum of integer squares
vanishes exactly when all rows vanish.  Exact expansion and collection
give the reported quartics.

The 2,560-variable sparse alternative retains semantic names for the Pell
outputs rather than making all of (10)'s definitions affine substitutions.
It has fewer collected monomials.  Its orientation, short-beta source,
CRT block and radius bound are otherwise the same.  The intermediate
shared-offset construction is also justified: finitely many signed
quotients can be written $v_i-h$ with one common natural $h$, chosen after
all other witnesses.  It is superseded by the forward form for the final
natural variants, except that one final CRT congruence still uses two
natural quotients.

For the 4,087-variable integer companion, congruence quotients themselves
need not be nonnegative.  There are exactly 197 unrestricted quotient
coordinates: four per positive-power graph and one final CRT quotient.
In (10), $S,R_0$ are retained as separate natural semantic coordinates,
with their defining equalities, while $q_S,q_R$ may be signed.  The other
parametric definitions are unchanged.  In particular $B_0>1$ still follows
from $B_0=1+4Y(q_4+1)$ with natural $q_4$; it is not incorrectly inferred
from signed $q_B$.  Soundness reads congruences directly and applies the
same indexed-Pell and power theorems.  Completeness uses the natural
witness construction already proved.  The final CRT quotient may be
signed because (6) is a divisibility condition.

Only 607 further nonnegativity assertions are necessary.  An exact sign
certificate infers the other 44 of the 651 nonquotient semantic
coordinates from polynomial equalities.  Each inference is of the form
$c x=P$ with $c>0$ and $P$ having nonnegative coefficients; every variable
occurring to an odd exponent in a monomial of $P$ must already be known
nonnegative.  Even powers need no sign premise.  The supplied application
order is an acyclic proof of all required signs.  Signed quotient gates
are not asserted to be nonnegative.

For each remaining seed $x$, append the quadratic residual
\[
 2x-a(a+1)-b(b+1)-c(c+1)=0. \tag{11}
\]
For arbitrary integers $a,b,c$, each consecutive product is nonnegative,
so (11) forces $x\ge0$.  Conversely every natural $x$ is a sum of three
triangular numbers by Gauss's theorem.  Equivalently, apply Legendre's
three-square theorem to $8x+3$: the three squares must all be odd modulo
8, and subtracting their three units yields (11).
The 607 seeds add 1,821 integer coordinates and 607 quadratic rows to
the 2,266-coordinate mixed-domain template.  This proves the counts
$2266+1821=4087$ and $1924+607=2531$, and preserves degree four.

\subsection{Source pins, reproduction, and the claim boundary}

The calibration files are pinned to
\texttt{VladimirReshetnikov/ProveIt} commit
\nolinkurl{08200eb10fde2db408c69dffd840d6920a569b1f};
the comparison source is the fixed-power-bound candidate at commit
\nolinkurl{7e533c33819ee1f703b69f2778a767cdc069940d}.
The exact theorem names, file paths and SHA256 values are listed in
\texttt{SOURCE\_PINS.json} and the inherited calibration inventory.
The all-forward Pell proof additionally pins the official Mathlib file
at commit \nolinkurl{4dc43842a11f257ffba2a611c943a133a65cb4b2},
SHA256 \nolinkurl{5c5756cdf2fa70fb77af9daa9c7167c65749ac616bde8f5fdfeed2cb958b3ed5}.
It is inspected as source text, not executed.

The builders are newly authored standard-library Python implementations.
No upstream executable was imported or run.  Their baseline mode
independently reconstructs all five published numerical baseline metrics.
The final release includes the exact sparse equations, every residual,
raw circuits and macro bindings, affine restoration maps, structural
radius certificates, integer sign certificates, source pins, and
independent finite-audit receipts.  Reproduction commands are in
\texttt{README.md} and the file whitelist is machine-readable.

These finite checks establish literal polynomial and certificate facts,
not consistency of ZFC and not the universal mathematical dependencies.
The logical chain is: pinned numerical refutation calibration; finite
single-conclusion proof flow and the carrier lemma; short-beta source;
factored Horn zero set; bounded CRT with factorial cancellation; radius bound; explicit arithmetic
identities and Pell graphs; quartic flattening; and, for unrestricted
integers, sign closure and the three-triangular-number theorem.  No
nonpolynomial operation remains in the final equation.

\subsection{Final source projection: eight columns and no private bound}
\label{sec:projected-direct}

The source projection in this subsection is explicitly credited to the
incoming \emph{ConZFC Direct Projection 7617} package, whose preprocessing
proof is preserved as \nolinkurl{sources/INCOMING_PROJECTION_PROOF.md}.
The new result here is its composition with the sharper carrier,
parametric positive-power graphs, dyadic majorants, quotient bound $x$,
and sign-normalized integer compilation.  We do not claim that eliminating
these three label columns was first discovered in this construction.

Each rule can be parameterized by just $a,b,c,d$:
\[
\begin{array}{c|c|c|c}
 \text{rule}&p&q&v\\ \hline
 \text{truth}&0&0&0\\
 \text{eight ordinary axioms}&0&0&f_j(a,b,c,d)\\
 \text{bundled axiom}&0&0&f_{\rm bundle}(a,b,c,d)\\
 \text{modus ponens}&a&I(a,b)&b\\
 \text{generalization}&b&0&A(a,b)
\end{array}
\]
The bundled row additionally requires
$((c-a)(c-b))^2=(d+1)^2$.  No other row has a guard.  These substitutions
are exactly the pinned C222 rules, including the independent
parameter and premise of generalization.

For a chosen branch $j$, substitute its labels $(p_j,q_j,v_j)$ into the
accumulator residual
\[
 L_j=S'(1+v_jM)-S(1+p_jM)(1+q_jM).
\]
Let $C_j=L_j^2$ except at the bundle, where
$C_j=L_j^2+g^2$.  The new local polynomial is
\[
 P_8=\prod_{j=0}^{11}C_j
 +\bigl(x-S-(1+(i+1)y)u\bigr)^2
 +\bigl(x-S'-(1+(i+2)y)u'\bigr)^2. \tag{12}
\]
Its zero set is exactly one valid branch-specific update together with
the two common beta reconstructions.  As all branch factors are
nonnegative, the product implements the finite disjunction without
selecting a rule by an extra coordinate.  This is existential projection
of the three deterministic labels, not an assertion that the polynomial
in (12) is algebraically identical to the eleven-coordinate source.

\paragraph{Padding chooses the parameter bound before the carrier.}
Given a finite proof-tree expansion, choose $N\ge1$ at least its length
and at least every schema parameter in the expansion.  For modus ponens,
this may include a very large formula code; no practical bound on the
resulting $N$ is claimed.  Pad to exactly $N$ rows using truth, whose
labels are $(0,0,0)$ and whose update leaves $S$ unchanged.  Thus all four
parameters may be bounded by $N$ without changing refutability.

Every branch label is a nonnegative-coefficient polynomial in those four
parameters.  Apply the structural exponent certificate to its literal
constructor DAG with radius $N+4$.  The largest certified exponent is
3,445,903, so the fixed polynomial
\[
 m=(N+4)^{2^{22}} \tag{13}
\]
bounds every label.  It is computed by exactly 22 squarings, not by a
variable-exponent graph and not by a duplicate evaluation of the bundled
axiom.  It satisfies $m\ge1$.  Choose $M=(m+1)^N$ and $M_0=M-1$.
The carrier proof already given applies verbatim to these labels.  The
choice order is proof, padded $N$, then $m,M$, then the accumulator
sequence, then $y,x$.  There is no self-referential bound.

The eight local coordinates are
\[
 (a,b,c,d,S,S',u,u'),\qquad
 \text{bounds }(N,N,N,N,y-1,y-1,x,x). \tag{14}
\]
The use of $x$ for the two quotient bounds is an additional simplification:
canonical beta reconstructions have
$x=S+Du$ with $S\ge0,D\ge1$, hence $u\le x$, and similarly for $u'$.
No new private bound $U$ is needed.  Conversely, the state bounds in (14)
make those reconstructions canonical, just as in the eleven-column proof.
The endpoints and parameterization $y=M_0+2+y_0$ are unchanged.
The four outer natural coordinates are $N_0,x,y_0,M_0$, with $N=N_0+1$.

\paragraph{A second dyadic radius proves the shared CRT bound.}
Split the actual signed circuit for (12) as $P_8=E_8-F_8$ and put
\[
 R=4+N+x+y+M_0,\qquad H=R^{2^{23}}. \tag{15}
\]
All inputs at the upper substitution of (14) are at most $R$; for the two
states we safely use $y$ instead of $y-1$.  The two root exponents from
the structural certificate are both 6,893,248, below $2^{23}$.
Therefore $H$ bounds both source evaluations and $N$.  It also dominates
every binomial cutoff $N+1,y,x+1$.  The same CRT theorem, with direct
factorial cancellation, applies.  Only 23 successive squarings compute
(15).  Notice that $m$ need not occur in this source radius: after label
projection the local source computes its labels from the four bounded
parameters and uses only $M_0$ as its carrier input.

The final eight-column arithmetic uses exactly 40 positive-power graphs,
11 binomial graphs and two factorial graphs.  The count is transparent:
24 powers for the eight coordinate binomials, five for each of the two
factorials, five other powers in the modulus product, and one carrier
power.  In this accounting the five other product powers are its two
explicit powers and the three powers of its extra binomial; the factorial
at $N$ has already been counted.  No arithmetic macro, variable-length
product or variable exponent remains after flattening.

\subsection{Existential sign normalization for the final integer equation}

The incoming signed-quartic package supplies the sign-orbit method
and its two pivots per positive-power graph; its proof is preserved as
\nolinkurl{sources/INCOMING_SIGN_NORMALIZATION_PROOF.md}.  The final
application here is checked against the new projected parametric system,
not assumed merely because it worked for a predecessor.

For a family of residuals $r_j$, a diagonal sign action
$x_i\mapsto(-1)^{\epsilon_i}x_i$ preserves the common zero set whenever
all monomials of each row have equal parity against $\epsilon$.
Then $r_j$ is changed to either $r_j$ or $-r_j$.
For each of the 40 power calls there are two independent actions:
\begin{itemize}
\item Flip its auxiliary $Z$ and the product wire $WZ$.  The only other
      use of that wire is through its square in the last norm.
\item Flip its Pell coordinate $U$ and its two unrestricted quotients
      $q_B,q_S$.  The products $Uq_B,Uq_S$ are unchanged, while the norm
      uses $U$ only through $U^2$.
\end{itemize}
The actions leave every other semantic coordinate unchanged.  They
commute, and do not change any retained sign guard.  At any integer zero,
apply a pivot's action exactly when that pivot is negative.  This produces
another zero in which all 80 pivots are nonnegative.  This is an
existence-preserving normalization, not a claim that the original zero
already had those signs.

The released normalization certificate specifies all flipped raw and
exported variables and the parity of every residual under each action.
The separate independent checker verifies these finite identities; the
argument above explains their universal meaning.  Gate values involving
signed quotients remain allowed to be negative.  No signs of $X,S,R_0$
are discarded by an unsupported symmetry argument.

The mixed-domain template has 1,937 variables and 1,660 quadratic rows.
It has 161 unrestricted congruence quotients: four per power call and
one final CRT quotient.  There are 529 other semantic coordinates.
Before normalization, exact polynomial sign closure requires 494 seeds
and infers 35 more signs.  The 80 pivots are among these seeds and can
instead be normalized.  Thus only
\[
 414+80+35=529
\]
semantic signs need to be guarded, normalized, or inferred, respectively.
The 414 three-triangular-number certificates add $3\cdot414=1242$
integer variables and 414 quadratic rows.  Consequently the final
integer counts are
\[
 1937+1242=3179,\qquad 1660+414=2074.
\]
The resulting ordinary quartic $Q_8$ has 23,157 monomials and coefficient
height 63,923,766.  The natural companion has 1,859 variables, 1,581 rows,
12,275 monomials and height 63,924,066.  Both heights are below $2^{26}$.
The literal files are \texttt{projected/expanded\_polynomial.json} and
\texttt{projected\_integer/expanded\_polynomial.json}.

\paragraph{Complete equivalence and comparisons.}
A zero of the integer polynomial first yields an integer zero of the
mixed template with the 414 guarded signs.  Normalize the 80 pivots and
replay the 35 sign implications.  All nonquotient arithmetic inputs are
then natural.  The parametric Pell certificates establish every required
power; formulas (7)--(9) establish every other arithmetic macro; the CRT
certificate recovers all eight-column rows; their parameter and label
bounds license the sharp carrier and hence a C222 refutation.  The pinned
calibration gives a ZFC refutation.  Conversely a finite refutation gives
the padded source in (12)--(14), all arithmetic certificates, natural
power witnesses, and the 414 triangular representations, hence an integer
zero.  Therefore
\[
 \operatorname{Con}(\mathrm{ZFC})
 \quad\Longleftrightarrow\quad
 \neg\exists z\in\mathbb Z^{3179}\;Q_8(z)=0.
\]
All steps are ordinary mathematics formalizable in ZFC.  Finite artifact
checks do not replace the stated unbounded arithmetic and calibration
theorems, and no new proof-assistant certification or proof of consistency
is claimed.

The late incoming direct integer alternatives have 3,776 variables,
20,194 monomials and height 47,959; and 3,671 variables, 26,271 monomials
and height 3,835,369,234.  The new 3,179-variable equation improves the
latter on variables, monomials and coefficient height, but does not
improve the former's monomial count or its much smaller coefficients.
It should not be described as dominating every current direct equation.
The 2,560-natural-variable sparse companion likewise remains relevant
because its 11,074 monomials are fewer than the projected natural
companion's 12,275.  These are explicit tradeoffs, not an absolute
minimality claim.

\section{Measured alternative certificate architectures}
\subsection{Negative CRT moduli as a valid alternative}

Suppose the row condition is $\prod_j C_j=0$ and $R_0=0$, where all
$C_j,R_0$ are sums of squares. Bound each $C_j$ and $R_0$ separately
by $K$, with $K\ge4,N$ and all range cutoffs. Put
\[
 T=K!,\quad D_i=(i+1)T-1,\quad B=\prod_{i<N}D_i.
\]
The $D_i$ are pairwise coprime by the usual difference argument, and each
is $3\pmod4$. Hence each has a prime divisor $p\equiv3\pmod4$;
since $p\nmid K!$, necessarily $p>K$.
Using the usual binomial range certificates, the index congruence
$B\mid(C+1)T-1$, and
\[
 \left(\prod_jC_j(C,A)\right)^2+R_0(C,A)^2\equiv0\pmod B
\]
recovers every row. Indeed, over such a prime, a sum of two squares is
zero only when both terms are zero. Primality then selects a vanishing
$C_j$, and the separate bounds turn the selected branch and $R_0$ into
exact zeros. Ordinary CRT proves completeness.

No variable-length product remains in the output. With
$Q=(1+(N+1)T)^N+2$ one has
\[
 B=\left(T^N N!\binom{Q+N}{N}\right)\bmod(TQ+1).
\]
The numerator is $\prod_{j=1}^N(TQ+jT)$, which reduces to
$\prod_{j=1}^N(jT-1)$, and $B<Q<TQ+1$. For $N=0$ the product is
$1$, the congruences modulo $1$ are vacuous, and the same identity holds.
The rounded-radius negative-modulus benchmark has 2678 natural variables,
six more than its matched positive-modulus comparison. It is a proved
alternative, not the recommended minimum.

\subsection{Factorial separation without a coefficient height}

If $M\ge1$ and $\max(m-1,0)!\mid M$, then the numbers
$1+jM$, $1\le j\le m$, are pairwise coprime. A common divisor for
labels $j<k$ is coprime to $M$ and divides $k-j$, while $k-j$ divides
$M$. Therefore it is $1$. Prime valuations now recover every positive
label multiplicity from a product equality, for arbitrary proof length.
The factors themselves need not be prime. Label zero contributes the
unit factor and remains invisible. For $m=0$ the target factor $1+M>1$
makes the equality impossible; for $m=1$ any positive $M$ works.

Thus $M=m!$ is a correct replacement carrier when $m\ge1$. It is
chosen before the obligation stream, never from its later CRT bound.
The matched native-graph benchmark costs 2872 instead of 2672 variables:
one extra factorial, one extra binomial and four extra power graphs.
A completely fixed $M$ without a label restriction is unsound, as shown by
$1+(M+2)M=(1+M)^2$ and the underivable single rule $(1,0)\to M+2$.

All arguments above are finite arithmetic theorems formalizable in ZFC.
They remain internal theorems in nonstandard models; they do not rely on
external standardness of a proposed certificate. The included finite tests
and literal polynomial replays verify implementation fidelity, not a new
kernel proof of those all-input statements.

\section{Coprime scale carriers and fixed prime separation}
\label{sec:r13-radicals}

A low-degree scale and fixed prime separation can be combined without
adding a variable. The resulting complete integer equations are
\begin{center}
\begin{tabular}{lrrr}\toprule
Construction & Variables & Exact degree & Binary operations\\\midrule
Normalized five-root norm &13&1072&192\\
Ordinary four-root norm &14&378&145\\
Ordinary three-root norm &15&172&123\\
Centered three-root norm &15&170&134\\
Centered two-root norm &16&98&138\\\bottomrule
\end{tabular}
\end{center}
All variables range over unrestricted integers. These are explicit
cost--degree tradeoffs, not global minima. In particular, the first row
improves the incoming $13/1088/195$ point; the second substantially lowers
both degree and evaluation cost from the inspected $14/652/163$ point.
The $16$-variable row lowers degree but increases cost relative to
$16/108/132$.

The source is the literal signed CRT parent
\texttt{degree88\_17.json}, with its inherited ordinary native compiler
and K215 calibration. We restore each of its five residuals below.
This is not a new Lean/Rocq certificate, and it proves neither consistency
nor a historical or global optimality claim.

\subsection{Complete common prefix}
The base thirteen coordinates are
\[
 j,C,z,h,r,v,w,T,c,n,t,p_2,p_3.
\]
The fixed six numerals are the coherent source-34 tuple
$\beta,K_0,M_C,M_F$, with $M_F$ already shifted, and
$A_{\rm cone}=2^{1113}$, $B_{\rm cone}=2^{666104}$.
In particular $\beta=2^{1064921}-1$. All named registers below are
polynomial abbreviations, not extra quantified variables:
\begin{align*}
J&=6j+1,&q&=\beta J+1,&W&=12w+4,\\
X&=qW,&Y&=qv,&E&=XY,\\
a&=E+Y,&H&=4a+3,&\Delta&=a^2+H,\\
F&=(K_0+W)C-z(q-1),&&&\\
R&=(q(q-F)-C)(q^2-1)+(M_C+qM_F)J,&&&\\
k&=R+1+hE,&\eta&=c-Yk,&\zeta&=k-\eta,\\
u&=EYk,&A_\tau&=u(u+k)+1,&A_m&=1+\Delta c^2,\\
b_0&=X+ac,&I&=(2r+1)c^2,&M&=2q^2,\\
U&=MH,&e&=v-q^2.&&
\end{align*}
The determinant and guard radicand are
\begin{align*}
d&=(A_{\rm cone}F)(B_{\rm cone}C)
 -(q-A_{\rm cone}F-B_{\rm cone}C)^2,\\
A_g&=d\eta\zeta-(t^2-d-\eta\zeta+1)^2
     -1-p_2^2-p_3(p_3+1).
\end{align*}
For the ordinary variants, add the coordinate $f$ and set
\[
 Q=I^2,\qquad S=Q-\Delta(f^2-1),\qquad
 V=fcT-f^2R-c.
\]
For the thirteen-variable normalized variant instead use
\[
 A_f=1+\Delta I^2,\qquad Q=\Delta^2I^2,\qquad
 V=A_f(cT-R)-c.
\]
In both cases put $\delta=Q-1$, $C_y=QV^2-1$. The rational expression
$A_y=C_y/\delta$ is proof notation only, never a circuit division.
Where the coordinate $\gamma$ is retained, write
$D=(b_0+H\gamma)^2-A_m$. Where $p_1$ is retained, write
$G=A_g-p_1^2$.

\subsection{Five literal norm specifications}
Let
\[
 \mathcal F_m(x;a_1,\ldots,a_m)
 =\prod_{\varepsilon\in\{-1,1\}^m}
   \left(x+\sum_i\varepsilon_i\sqrt{a_i}\right).
\]
Sign invariance proves that this is an integer polynomial. The following
expressions are cleared symbolically; their JSON circuits use only binary
addition, subtraction and multiplication.

For thirteen variables take $Z=U(n+V)+Mb_0+He$ and
\[
 N_{13}=\delta^{16}\mathcal F_5\left(
 Z;81U^2A_f,9U^2A_\tau,M^2A_m,729U^2A_g,
 \frac{U^2C_y}{\delta}\right).
\]
The equation is $N_{13}=0$.

For fourteen variables, add $f$, retain the same $Z$, and use
\[
 N_{14}=\delta^8\mathcal F_4\left(
 Z;9U^2A_\tau,81U^2A_g,M^2A_m,
 \frac{U^2C_y}{\delta}\right),\qquad
 P_{14}=(1+N_{14})(1+S)-1.
\]
For fifteen variables, add $f,\gamma$, and use
\[
 N_{15}=\delta^4\mathcal F_3\left(
 M(n+V)+e;9M^2A_\tau,81M^2A_g,
 \frac{M^2C_y}{\delta}\right),
 \qquad P_{15}=(1+N_{15})(1+S)(1+D)-1.
\]

The centered variants use $t_0=2u+k$, $\ell=k^2-4$.
For fifteen variables add $f,p_1$ and set
\begin{align*}
 Z_c&=U(n+3t_0+2V)+2Mb_0+He,\\
 N_{15c}&=\delta^4\mathcal F_3\left(
 Z_c;36U^2A_\tau,4M^2A_m,\frac{4U^2C_y}{\delta}\right),\\
 P_{15c}&=(1+N_{15c})(1+S)(1+4G)-1.
\end{align*}
For sixteen variables add $f,\gamma,p_1$, put $x=Mn+e$, and define
\begin{align*}
 j_0&=\delta\{x^2+M(6t_0+4V)x+9M^2\ell\}
       -4M^2(V^2-1),\\
 h_0&=12\delta t_0VM^2,\\
 N_{16}&=j_0(j_0+2h_0)
  +144\delta M^4\{4A_\tau(1-V^2)+\delta\ell V^2\},\\
 P_{16}&=(1+N_{16})(1+S)(1+D)(1+4G)-1.
\end{align*}
The exact identity, valid also at $\delta=0$, is
\[
 N_{16}=\delta^2\mathcal F_2\left(
 M(n+3t_0+2V)+e;36M^2A_\tau,
 \frac{4M^2C_y}{\delta}\right)
\]
after symbolic clearing. An independent integer-coefficient check compares
all $29$ nonzero terms in the centered identity.

\subsection{Unit products recover the rows before denominator use}
For integers, if residuals $R_i$ never equal $-2$, then
\[
 (1+N)\prod_i(1+R_i)=1
 \quad\Longleftrightarrow\quad N=R_1=\cdots=R_s=0.
\]
Indeed every factor of the integer unit $1$ is $\pm1$. Each $R_i$ factor
is $+1$, forcing the one unrestricted norm factor to be $+1$ too.
There is no missing assumption on the norm modulo four.

Here $a$ is even and $\Delta\equiv3\pmod4$. If $f$ is odd then
$S\equiv I^2\pmod4$; if $f$ is even then $S\equiv I^2-1\pmod4$.
Thus $S$ is never $2$ modulo four. The same is true of $D$: for even $c$
its residues are $0,3$, and for odd $c$ they are $0,1$.
Finally $4G\equiv0\pmod4$. All the stated finalizers therefore recover
their individual residuals, in particular $S=0$, before interpreting a
rational norm.

Since $a\in6\mathbb Z$, $\Delta=(a+1)(a+3)\ge3$.
In ordinary variants $S=0$ excludes $Q=1$, as it would give
$\Delta(f^2-1)=1$. In the normalized variant $Q=\Delta^2I^2$ excludes
$Q=1$ directly. Hence $\delta\ne0$ at every zero in question.

\subsection{Prime separation before guard recovery}
The fixed residues give
\[
 q\equiv X\equiv2\pmod3,\quad a\equiv H\equiv\Delta\equiv0\pmod3,
 \quad A_\tau\equiv A_m\equiv1\pmod3.
\]
To check $A_\tau$, write $A_\tau=1+b(b+1)k^2$ with
$b=XY^2\equiv0$ or $2$ modulo three.
Also $H$ is nonzero and odd. Under $S=0$, the square $Q=I^2$ is divisible
by three, hence by nine. The normalized $Q$ is already divisible by nine.
Thus $A_y$ is a three-adic unit of residue one.

The weighted roots used by the ordinary fourteen-variable norm have
squared orders
\[
 -2v_3(H),\quad0,\quad2,\quad4+v_3(A_g)
\]
for $d_0/H,y,3\tau,9p_1$ respectively. The normalized construction
inserts $9f$ of squared order four and uses $27p_1$ of squared order
at least six. A zero guard root is harmless.

Here is the precise separation argument. In any subset of at least two
nonzero roots, the smallest squared order is unique and even. Divide
out its power of three and multiply the putative zero sum over all signs
of the remaining roots. The product is a polynomial in the rational
squares and is a nonzero unit modulo three, a contradiction.
Subtracting a Galois conjugate of a rational signed root sum would give
exactly such a forbidden signed subsum. Every root is therefore rational.
This argument allows negative radicands; their rationality subsequently
rules that case out. No positivity guard is assumed beforehand.

Roots of integer squares are integral. For the auxiliary root, $Q=0$
gives $y=\pm1$. For $Q>1$, rational $y$ in
$QV^2-(Q-1)y^2=1$ is integral by adjacent-Pell descent. In lowest terms
$y=A/m$ gives $m^2\mid Q-1$. The transformation
\[
 V'=(2Q-1)V-2(Q-1)y,\qquad
 y'=(2Q-1)y-2QV
\]
preserves the norm and the denominator, and after taking positive signs
satisfies $0<V'<V$ when $V>1$. The endpoint $|V|=1$ forces $m=1$.
Thus the case $c=0$, and consequently $Q=0$, has not been silently excluded.

\subsection{Coprime carriers and exact source restoration}
Because $Y=qv$, we have $H\equiv3\pmod q$. Since $q\equiv2\pmod3$
and $H$ is odd,
\[
 \gcd(H,2q^2)=1.
\]
After root recovery the ordinary fourteen-variable factor is
\[
 U(n+V)+Mb_0+He=Md_0+U(3\tau+y+9p_1).
\]
Reduction modulo $M$ and $H$ gives respectively
\[
 v=q^2(2v_{\rm old}+1),\qquad
 \gamma=\frac{d_0-b_0}{H}\in\mathbb Z.
\]
The normalized case adds $9f$ and replaces $9p_1$ by $27p_1$;
the same argument applies. The centered main carrier instead has
$2M(d_0-b_0)$ modulo $H$, still invertible because $H$ is odd.
The variants retaining $\gamma$ need only the first congruence.
In all cases $Y=(2v_{\rm old}+1)q^3$ exactly, including its odd quotient.
The removed $v^2$ guard multiplier is unnecessary: $v\ne0$ follows
now from the recovered congruence, not from an assumed inequality.

The integral guard root restores $G=0$, which implies $d>0$ and
$\eta\zeta>0$ by the standard joint-positivity identity. For completeness,
if both factors were negative their magnitudes $b,e>0$ would require
$(t^2+b+e+1)^2<be$, which is impossible. Thus $c=0$ is impossible too,
since it would give $\eta\zeta=-Y(Y+1)k^2\le0$.

For the ordinary variants set old $J_3=2j$, old scale $v_{\rm old}$
as above, and restore the two centered source coordinates by
\[
 \tau_{\rm old}=2\tau-(2u+k),\qquad y_{\rm gap}=y-V.
\]
The first and auxiliary parent residuals become exactly
\begin{align*}
 \tau_{\rm old}^2+2\tau_{\rm old}(2u+k)+k^2-4
   &=4(\tau^2-A_\tau),\\
 V^2+(1-Q)y_{\rm gap}(2V+y_{\rm gap})-1
   &=QV^2-(Q-1)y^2-1.
\end{align*}
The other three rows are the main norm, $S$ and $G$.
Every remaining parent abbreviation agrees literally. This restores the
five-row signed CRT source directly, without inventing a new rank lemma.
For normalized thirteen variables instead set
\[
 r_{\rm old}=\frac{\Delta(2r+1)-1}{2},\qquad
 T_{\rm old}=fT.
\]
The first quotient is integral because $\Delta$ is odd. Then
$I_{\rm old}=\Delta I$, the parent's $V$ agrees exactly, and its strong
row follows from $f^2=A_f$.

Conversely, the inherited canonical refutation witnesses have odd
$J\equiv1\pmod3$, so $J=6j+1$. Put
$v=q^2(2v_{\rm old}+1)$ and retain their guard witnesses.
For example the ordinary fourteen-variable construction chooses
\[
 n=\gamma+3\tau+y+9p_1-V-e/M.
\]
The centered fifteen-variable construction chooses
\[
 n=6\tau+2y+2\gamma-3t_0-2V-e/M.
\]
The other formulas omit the coordinates not fused by their norms.
All values are integers and make one conjugate factor vanish.
The normalized thirteen-variable case uses the same canonical odd strong
normalization and fresh auxiliary reconstruction proved in the incoming
prime-separation package. Its ordinary calibration boundary is unchanged.

\subsection{Universal degree bounds and actual coefficient certificates}
Fix the six ports. The exact propagated prefix bounds are
\[
\begin{array}{c|rrrrrrrrrr}
\text{register}&H&k&u&A_\tau&A_g&Q&V&M&U&S,D\\\hline
\text{degree}&4&5&11&22&28&6&7&2&6&10
\end{array}
\]
in the ordinary case, with $\deg G=28$. The normalized case instead has
$\deg A_f=14$, $\deg Q=22$ and $\deg V=19$.
Direct degree propagation through the exact DAGs proves the final upper
bounds $378,172,98$ for ordinary fourteen, ordinary fifteen and centered
sixteen variables. Two cancellation-sensitive cases require more.

\paragraph{Thirteen variables.}
Pair the two signs of the auxiliary root first. Write
\[
 Z=UV+s,\qquad s=Un+Mb_0+H(v-q^2),
\]
and, for each of the sixteen sign choices, put
\[
 B=s\mathbin{\pm}9Uf\mathbin{\pm}3U\tau
     \mathbin{\pm}Md_0\mathbin{\pm}27Up_1.
\]
Assign formal root weights $7,11,5,14$ to $f,\tau,d_0,p_1$.
Each squared radicand has degree at most twice its root weight, so
reduction in this formal root algebra never increases weighted degree.
Thus $B$ has weight at most $20$. The paired factor is exactly
\begin{align*}
 \delta\{(UV+B)^2-U^2C_y/\delta\}
 &=\delta B^2+2\delta UVB-U^2(V^2-1).
\end{align*}
The identity uses $C_y=(\delta+1)V^2-1$. Its three terms have bounds
$62,67,50$. The sixteen-factor product therefore has degree at most
$16\cdot67=1072$, and is exactly the cleared five-root norm.
This argument is polynomial algebra, including all exceptional inputs.

\paragraph{Centered fifteen variables.}
Use the exact centered two-root identity with its scale $M$ replaced by
$U$, and its carrier argument replaced by
\[
 x=Un+2Mb_0+H(v-q^2)\mathbin{\pm}2Md_0.
\]
With $d_0$ of weight five, $x$ has weight at most seven. Define
\begin{align*}
 j&=\delta\{x^2+U(6t_0+4V)x+9U^2\ell\}
      -4U^2(V^2-1),\\
 h&=12\delta t_0VU^2.
\end{align*}
Their weights are at most $30$ and $36$. Every term of
\[
 j^2+2hj+576\delta U^4A_\tau(1-V^2)
       +144\delta^2U^4\ell V^2
\]
has weight at most $66$. The exact identity uses only
$t_0^2=4A_\tau+\ell$, not a Pell zero. Taking the two main-root
conjugates gives norm degree at most $132$; the strong and guard
finalizers add $10+28$, yielding $170$.

Independent modular-line certificates modulo $65521$ and $65519$, using
the actual six numeral byte strings, attain every stated bound with a
nonzero leading coefficient. They are rigorous lower certificates
combined with the preceding universal upper proofs.
All literal DAGs pass topology, variable closure, liveness and coefficient
hash checks. An independent exact sparse-algebra audit verifies every
prefix, carrier, radicand, norm and finalizer, using all $20\,349$
coefficients of the generic five-root norm; deliberate corruptions are
rejected. Signed modular diagnostics additionally include $c=0$, $v=0$
and their intersection. These algebraic checks do not establish an
unbounded compiler theorem: the named ordinary source calibration remains
an explicit mathematical dependency.

\section{Twelve integers at degree 2048}
\label{sec:twelve2048}
A late supplied interpolation theorem permits a further synthesis. The
result is a literal polynomial in twelve unrestricted integer variables,
of exact degree $2048$, evaluated in $233$ live binary operations
($148$ multiplications, $53$ additions and $32$ subtractions).
It improves the supplied twelve-variable $4832/248$ degree--cost point.
The supplied $16000/223$ point remains a cheaper, higher-degree alternative.
No fewest-variable or global optimality claim is made.

The new changes are a coprime homogeneous quintic carrier, a guard Pell
witness divisible by three, fixed prime separation in place of large
size weights, and a doubled auxiliary-root center. The full-residue
interpolation theorem and the direct normalized signed-kernel calibration
are imported from \texttt{ConZFC\_Twelve\_Integer\_Interpolation\_2026-10-24}.
Its exact proofs, provenance and notices are retained. In particular,
this is not a literal restoration of the earlier odd-multiplier source
chart, nor a new proof-assistant certificate.

\subsection{Complete polynomial specification}
The twelve coordinates are
\[
 j,C,z,h,r,v,w,T,c,n,t,s_0\in\mathbb Z.
\]
Use the same complete six fixed source-34/source-71 numerals as before.
Define
\begin{align*}
 J&=6j+1,&q&=\beta J+1,&W&=12w+4,\\
 X&=qW,&Y&=qv,&E&=XY,\\
 a&=E+Y,&H&=4a+3,&\Delta&=a^2+H,\\
 F&=(K_0+W)C-z(q-1),&&&\\
 R&=(q(q-F)-C)(q^2-1)+(M_C+qM_F)J,&&&\\
 k&=R+1+hE,&\eta&=c-Yk,&\zeta&=k-\eta.
\end{align*}
Put
\begin{align*}
 d_g&=(A_{\rm cone}F)(B_{\rm cone}C)
          -(q-A_{\rm cone}F-B_{\rm cone}C)^2,\\
 P&=d_g\eta\zeta-(t^2-d_g-\eta\zeta+1)^2,\\
 s&=3s_0,&D_g&=(3P-4)s^2+1,&I&=src^2,\\
 D_f&=1+\Delta I^2,&Q&=\Delta^2 I^2,&\delta&=Q-1,\\
 V&=D_f(c-R+4IT)-c,&C_y&=QV^2-1,&D_a&=\delta C_y,\\
 \ell&=EYk,&D_\tau&=1+\ell(\ell+k),&D_m&=1+\Delta c^2,\\
 b_0&=X+ac,&M&=2q^2,&L&=MH,\\
 e&=v-q^2,&\widetilde b&=Mb_0+He,&
 u&=LT-\widetilde b,\\
 \widetilde D&=M^2D_m.&&&
\end{align*}
For the fixed intersective quintic
\[
 \pi(x)=(x^2+31)(x^3+x+1),\qquad
 \Phi_L(x)=L^5\pi(x/L),
\]
the following integer polynomials are its exact quadratic remainder:
\begin{align*}
 A={}&u^5+10u^3\widetilde D+5u\widetilde D^2
       +32L^2(u^3+3u\widetilde D)\\
    &+L^3(u^2+\widetilde D)+31L^4u+31L^5,\\
 B={}&5u^4+10u^2\widetilde D+\widetilde D^2
       +32L^2(3u^2+\widetilde D)+2L^3u+31L^4.
\end{align*}
Thus $A+B\widetilde d=\Phi_L(u+\widetilde d)$ when
$\widetilde d^2=\widetilde D$. Finally set
\[
 S=IL^5,\qquad C_a=(2Q-1)V,\qquad Z=S(n+C_a)-A.
\]
With $\mathcal F_5$ denoting the ordinary $32$-sign norm polynomial,
the complete equation is
\[
 \boxed{\mathcal P_{12}
 =\mathcal F_5\bigl(Z;B^2\widetilde D,
 4S^2D_a,36S^2D_f,324S^2D_\tau,2916S^2D_g\bigr)=0.}
\]
There is no denominator clearing in this norm: all five radicands are
integer polynomials. The displayed homogeneous notation is expanded in
the actual DAG. No square root, division, variable exponent or additional
quantified register occurs in the exported equation.

\subsection{Exclude the entire zero multiplier branch first}
For all integer inputs $q$ is nonzero, even and $2$ modulo three;
$X\equiv2\pmod3$ and $a\in6\mathbb Z$. Hence
$\Delta=(a+1)(a+3)\ge3$, and $H$ is nonzero, odd and divisible by three.
Also $H\equiv3\pmod q$, so $\gcd(H,M)=1$ and $L\ne0$.
Thus $D_m\ge1$ and $\widetilde D>0$, including $v=0$.

The imported elementary quintic lemma says that $\pi$ has no zero in
a real extension of $\mathbb Q$ of degree at most two. The quadratic
factor has no real zero, while the cubic has degree three and is
irreducible by the rational-root test. Consequently
\[
 A+B\sqrt{\widetilde D}\ne0,
 \qquad A-B\sqrt{\widetilde D}\ne0.
\]
If $I=0$, then $S=0$, the last four norm radicands vanish, and the entire
polynomial specializes to
\[
 (A^2-B^2\widetilde D)^{16}\ne0.
\]
Every zero therefore has $I,S,s_0,r,c\ne0$. This is established before
dividing a factor relation by $S$. Since $Q=(\Delta I)^2$, at such a
zero $Q>1$ and $\delta\ne0$.

\subsection{Small fixed weights replace all size separators}
Modulo three the new carrier has
\[
 M\equiv2,\quad L\equiv0,\quad b_0\equiv2,
 \quad u\equiv2,\quad\widetilde D\equiv1.
\]
Thus $B\equiv5+10+1\equiv1\pmod3$. Also
\[
 D_f\equiv D_\tau\equiv D_a\equiv D_g\equiv1\pmod3.
\]
For $D_a$ use $Q\equiv0\pmod9$; for $D_g$ use $s=3s_0$.
For $D_\tau$, the factorization
$1+(XY^2)(XY^2+1)k^2$ has $XY^2\equiv0$ or $2$ modulo three.

Divide a vanishing factor by the already nonzero $S$.
The five weighted roots
\[
 \frac{BMd}{S},\quad 2\widetilde y,\quad6f,
 \quad18\tau,\quad54g
\]
have squared three-adic orders respectively
\[
 -2v_3(S),\quad0,\quad2,\quad4,\quad6,
\]
where $d^2=D_m$, $\widetilde y^2=D_a$, $f^2=D_f$,
$\tau^2=D_\tau$, $g^2=D_g$.
Each subset of at least two has a unique even minimum. Normalize its
smallest root by the corresponding power of three and multiply a putative
zero sum over all signs of the other roots. The product is a polynomial
in the rational squared roots and is a nonzero unit modulo three, a
contradiction. Subtracting a Galois conjugate of the rational factor
relation would give just such a forbidden signed subsum. All roots are
therefore rational. The nonzero rational multiplier $BM/S$ causes no
loss: $d$ is rational too. Roots of integer radicands are integral.
No size estimate, radical positivity assumption or bound on $R$ was
used for this extraction. Global monotonicity of the chosen quintic,
although true, is not needed by this new proof: $B\equiv1\pmod3$
already supplies its nonvanishing.

Now $D_g=g^2\ge0$. Since $s\ne0$, the alternative $P\le0$ would give
$D_g=(3P-4)s^2+1\le-3$. Thus $P>0$, and the joint-positivity identity
gives $d_g>0$ and $\eta\zeta>0$.
The rational value $y=\widetilde y/\delta$ satisfies
\[
 QV^2-(Q-1)y^2=1.
\]
Adjacent-coefficient Pell descent makes $y$ integral.

\subsection{One monic carrier restores both congruences}
After absorbing the factor signs into the integer roots, the norm relation
gives
\[
 S\mid A+BMd=\Phi_L(Md-\widetilde b+LT).
\]
In particular $L^5$ divides this homogeneous monic quintic.
For a prime with $p^a\Vert L$, if its argument had order $b<a$,
the monic fifth-power term would have order $5b$, strictly below every
other term. That contradicts divisibility by $p^{5a}$. Therefore
\[
 L\mid M(d-b_0)-He.
\]
Reduction modulo the coprime factors $M,H$ gives
\[
 M\mid e,\qquad H\mid d-b_0,
 \qquad v=q^2(2v_{\rm old}+1),\qquad
 \gamma=(d-b_0)/H\in\mathbb Z.
\]
Thus $Y=(2v_{\rm old}+1)q^3$ is recovered with its full odd quotient.
Moreover the quintic argument divided by $L$ is exactly
\[
 T+\gamma-e/M.
\]
This proves both congruences without adding a variable or presupposing
that $v$ is nonzero.

The late source's direct normalized signed-kernel proof now applies.
The physical coordinates satisfy its first, main and normalized strong
Pell equations, strict guards and native transport/index equations.
Its two auxiliary quotients are restored by
\begin{align*}
 o&=f(c-R+4IT),\\
 j_{\rm aux}&=\Delta s^2r^2c^3(c-R)+4srcD_fT,
\end{align*}
so $V=fo-c=cj_{\rm aux}-R$. The normalized rank argument uses
$c^2\mid I$ and does not require $I/c^2$ to be odd. Accordingly, we do
not assert that the old coordinate $(\Delta sr-1)/2$ is integral.
The inherited signed step-down, residue orientation and native compiler
then produce a finite syntactic refutation, with their ordinary theorem
status unchanged.

\subsection{Completeness and the exceptional guard value}
Start with the late source's canonical native witnesses and set
$v=q^2(2v_{\rm old}+1)$. Choose its nearest-square coordinate $t$ so
$P>0$. A crucial exception must be excluded before forcing $3\mid s$:
if $P=1$, the guard equation would be $g^2+s^2=1$ and no nonzero
multiple of three could serve as $s$.

Here that exception is impossible. Both cone products and $q$ are even,
so $d_g\equiv0\pmod4$. Therefore
\[
 P\equiv-(t^2-d_g-\eta\zeta+1)^2\equiv0\text{ or }3\pmod4.
\]
Positive $P$ is at least three. Thus $D=3P-4$ is positive and
$2$ modulo three, hence nonsquare. Take a positive Pell pair
$g_1^2-Ds_1^2=1$. Modulo three, $g_1^2+s_1^2=1$, so
$3\mid g_1s_1$. Squaring the unit gives
\[
 g=g_1^2+Ds_1^2,\qquad s=2g_1s_1>0,
 \qquad 3\mid s.
\]
Set $s_0=s/3$. No outer coordinate has changed.

Invoke the imported full-residue interpolation lemma with $N=s$.
It supplies a single normalized Pell pair $I,f$, with $4sc^2\mid I$,
\emph{before} a residue is requested, and then realizes every class
of $T$ modulo this same $I$ in
\[
 V=f^2(c-R+4IT)-c,\qquad
 (\Delta I)^2V^2-((\Delta I)^2-1)y^2=1.
\]
Put $r=I/(sc^2)$. The lemma's hypotheses include even $a+2$,
$R\equiv c\equiv3\pmod4$, all supplied by the canonical source.
Its proof uses the corrected period $4n_0f$ modulo $f^2$ and the
Lucas valuation bound $\gcd(4n_0f,4I^2)\mid4I$; mere infinitude of
Pell solutions would not suffice.

Choose an intersective residue for
$\pi(T+\gamma-e/M)$ modulo $I$. Interpolation supplies $T,y$ in that
class, leaving $I,f$ fixed. Hence
\[
 A+BMd=L^5\pi(T+\gamma-e/M)
\]
is divisible by $S=IL^5$. With $\widetilde y=\delta y$, choose
\[
 n=\frac{A+BMd}{S}+2\widetilde y+6f+18\tau+54g-(2Q-1)V.
\]
This is integral and makes one norm factor vanish. The witness order is
native data, then $t,s,s_0$, then $I,f,r$, then the intersective residue
and $T,y$, and finally $n$. There is no circular residue choice.

\subsection{Exact degree 2048}
The relevant polynomial degrees are
\[
\begin{array}{c|rrrrrrrrrr}
\text{register}&H&\Delta&P&I&D_f&Q&V&D_a&D_g&L\\\hline
\text{degree}&4&8&28&4&16&24&21&90&30&6
\end{array}
\]
and
\[
 \deg A=35,\quad\deg B=28,\quad\deg S=34,
 \quad\deg C_a=45,\quad\deg D_\tau=22,
 \quad\deg\widetilde D=14.
\]
Pair the two auxiliary signs first. Let $B_*$ denote $Sn-A$ plus
the four signed non-auxiliary radical terms. Their weights are
$35,42,45,49$ for main, strong, first and guard, so $B_*$ has weight
at most $49$. The paired factor is exactly
\[
 B_*^2+2SC_aB_*+S^2(V^2+4Q-4),
\]
because the polynomial identity
\[
 ((2Q-1)V)^2-4(Q-1)(QV^2-1)=V^2+4Q-4
\]
holds universally. Its term bounds are $98,128,110$. Taking sixteen
sign choices proves degree at most $16\cdot128=2048$.
The unique leading paired term is
$2SC_a(\mathbin{\pm}54S\sqrt{D_g})$; the eight negative signs cancel.
Thus the full leading form is
\[
 108^{16}S_0^{32}(C_a)_0^{16}(D_g)_0^8\ne0.
\]
More explicitly, the coefficient of
\[
 j^{928}w^{288}v^{320}h^{32}s_0^{128}r^{112}c^{224}T^{16}
\]
is
\[
 \boxed{2^{2064}3^{1400}\beta^{928}>0.}
\]
The exponent sum is $2048$. The $h^{32}$ term uniquely selects
$ha_0$ throughout the top part of $k^{32}$; $T^{16}$ then uniquely
selects $4I_0T$ throughout $(4I_0T-R_0)^{16}$. The remaining factors
are monomials, so no other fixed coefficient can cancel this term.

Independent actual-coefficient line certificates modulo $65521$ and
$65519$ also attain degree $2048$. The complete $233$-operation DAG
has no dead gates or hidden input leaves. Independent arithmetic checks
use a separate binomial quintic remainder and full $32$-sign product,
including every intersection of the exceptional slices $s_0=0,r=0,c=0$.
These finite checks do not replace the unbounded interpolation or
native calibration proofs. The resulting consistency equivalence is an
ordinary mathematical theorem composed with those explicitly retained
sources, not a proof of consistency or a completed object-language
ZFC derivation.

\section{A sharp limit for final coordinate centering}
The ten-integer equations now have the following literal circuit counts.
The comparison uses the current pinned short-guard family, rather than an
older incoming package. These changes preserve the corresponding
polynomials exactly; they do not improve the variable counts or degrees.
\begin{center}
\begin{tabular}{lrrr}
Variant & Prior operations & New operations & Exact degree\\\hline
Polar &471&467&1,245,324\\
Linear &475&470&1,179,788\\
Cubic &478&473&1,114,252\\
Six term &512&507&1,080,888
\end{tabular}
\end{center}
Every live binary addition, subtraction and multiplication costs one,
including squares and scalar products. Fixed integer leaves and reuse cost
zero. The full graphs, actual coefficients and independent checks are in
\texttt{norm/}; no intermediate register adds an unknown.

The more substantive result is a limitation of this entire centering route.
With the short-guard prefix fixed, no polynomial translation of the final
coordinate alone can produce degree below $1{,}080{,}888$. This includes
arbitrary real polynomial translations, not only the integral corrections
previously considered. It is not a lower bound for different guards, norm
carriers, source encodings, or changes of other coordinates.

\subsection{The fixed norm and its leading forms}
Retain the ten coordinates $(t,g_0,c_0,k_0,l,w,x,y,z,m)$ and the source
registers from the supplied short-guard parent. In particular,
\[
 \Delta=g(\mathrm{cap}-g)K^2-\mathrm{cap}^2(C-KL)^2,
 \qquad a_0=r+(9\Delta-3)(rz)^2,\qquad a_1=DFI.
\]
Set $Z=\Sigma m-\Theta$. The unpolarized final norm is
\[
 N=(Z^2-\Sigma^2(a_0+a_1))^2-4\Sigma^4a_0a_1.
\]
For a nonzero polynomial $P$, write $P_{\rm top}$ for its highest
homogeneous part. The relevant degrees are
\[
\begin{array}{c|rrrrrrr}
 P&g&\mathrm{cap}&K&L&r&a_0&a_1\\\hline
 \deg P&2&233&28548&4220&114189&294382&458818.
\end{array}
\]
The first term of $\Delta$ has degree $57331$, while its negative square
term has degree $66002$. Therefore, with
\[
 U=3\,\mathrm{cap}_{\rm top}K_{\rm top}L_{\rm top}r_{\rm top}z,
 \qquad V=(2AEGH)_{\rm top},
\]
we have
\[
 (a_0)_{\rm top}=-U^2,\qquad (a_1)_{\rm top}=V^2.
\]
The second identity also follows from the exact defect identity:
$\deg(a_1-(2AEGH)^2)=426050<458818$. Both $U$ and $V$ are nonzero.
Since $\deg\Sigma=81922$, the fixed cross term has degree
\[
 d=4(81922)+294382+458818=1080888
\]
and highest homogeneous part
\[
 (-4\Sigma^4a_0a_1)_{\rm top}=(2\Sigma_{\rm top}^{2}UV)^2.
\]
This is a nonzero square, not merely a coefficient whose sign was sampled.

\subsection{The lower bound for every polynomial translation}
Let $R$ be any real polynomial in the first nine coordinates, and define
\[
 B_R=[\Sigma(m+R)-\Theta]^2-\Sigma^2(a_0+a_1).
\]
The translated norm is $B_R^2-4\Sigma^4a_0a_1$. If
$2\deg B_R>d$, its degree is $2\deg B_R>d$. If $2\deg B_R<d$, its
degree is $d$. At equality, its degree-$d$ part is
\[
 (B_{R,\rm top})^2+(2\Sigma_{\rm top}^{2}UV)^2.
\]
A sum of squares of real polynomials can vanish identically only when
both summands vanish identically. The second is nonzero, so cancellation
is impossible. Thus every such translated norm has degree at least $d$.
The six-term shift attains $d$, proving the restricted-family optimum.
The old odd-degree argument does not apply after the short-guard change;
this positive-leading-square argument replaces it.

Further degree progress in this ten-variable family must therefore change
more than the final polynomial translation: for example the guard, the norm
carrier, the prefix, or another coordinate. This conclusion does not rule
out any of those routes.

\subsection{The exact arithmetic savings}
Write $t_c=t_{c,0}-2$ in the unchanged source. Reuse
\[
 X_+=n_n(s_c+t_{c,0}-1)-s_c,
 \quad X=X_+-1,
 \quad K=X_++\mathrm{twice\_k}(P_{1/2}-1).
\]
This removes two unnecessary constant adjustments. Since $F=F_{-1}+1$,
use $G=hF_{-1}-F$ instead of $(h-1)F_{-1}-1$. Share
$S^3=SS^2$ and $S^5=S^2S^3$ between $\Sigma$ and $\Theta$.
For the linear and cubic shifts, share $GH$ by writing
$I=(GH)^2-H^2+1$.

For the six-term shift put $u=C^2$, $h=A/2$ and $E_0=u\xi$. Its compact
factorization uses
\[
 u^7\xi^4=u^3E_0^4,\qquad u\xi^2=E_0\xi.
\]
These save two multiplications without changing the existing shift.
Together the exact identities account for the counts in the table.

\subsection{Verification and mathematical boundary}
The independent checker imports neither the producer nor its circuit
helper. It checks exact local polynomial equalities for every retained
source register and the complete norm, all actual coefficient payloads and
recipes, topology, backward liveness, operation counts, the expanded radical
defects, and nonzero actual-coefficient degree certificates. It also performs
$360$ signed modular native replays per variant and rejects four final-sign
mutations. Normal and optimized Python runs agree.

The new algebra preserves the inherited ordinary mathematical equivalence
to finite ZFC refutability. Its compiler, Pell and interpolation interfaces
remain explicit dependencies. The result is not a consistency proof or a
new end-to-end proof-assistant certificate. No new low-degree integer chart
was established here: the incoming $23/10$ and $20/18$ points remain prior
work rather than contributions of this iteration.

\section{What the checks establish}
\label{sec:checks}
The release separates three different levels of evidence.
\begin{enumerate}
\item \textbf{Literal finite algebra.} Independent programs re-expand the
quadratic rows of the selected direct quartics and compare every collected
integer coefficient. They verify every retained variable, affine
restoration, source binding, Pell macro, radius certificate, and the
integer sign closure. Fixed norm payloads are authenticated and their
arithmetic graphs checked for topology, liveness and operation counts.
\item \textbf{Universal mathematical reductions.} The carrier bootstrap,
beta reconstruction, prime-lift CRT, forward Pell witnesses, parametric
elimination, unit-product finalizers, root separation and degree upper
bounds are proofs about arbitrary inputs. Finite samples are not offered
as substitutes for them. Exact modular leading coefficients, once an
upper bound is proved, are valid certificates that the degree is attained.
\item \textbf{Logical calibration.} The direct route retains the named
numerical C222 and ZF/ZFC refutation interfaces. The low-variable route
retains the native compiler and K215 interfaces. Their ordinary source
proofs are included or identified at their explicitly recorded source pins. A new complete
Lean or Rocq derivation is not claimed. The package does not search for,
or decide the existence of, a zero of a consistency polynomial.
\end{enumerate}

The exact audit programs use explicit failures rather than removable
Python assertions. The release driver invokes the selected suites both
normally and with Python optimization enabled. It also regenerates the
selected equations in temporary directories and compares their bytes,
not merely their numerical summaries. Source and proof hashes, payload
hashes, check outputs and timing receipts are recorded separately.
Checks that write receipts are run on temporary copies, so verification
does not silently invalidate the release manifest.

The archive README lists the exact suite inventory and supported command.
The runtime receipt records actual counts and elapsed times of the replay;
it is not a promise about speed on another machine. Python 3 and SymPy
are required for the complete suite. The direct, architecture and radical
checks themselves use standard-library Python. No network access,
repository checkout, Lean installation or earlier round's ZIP is required.
Rebuilding this PDF additionally requires a conventional LaTeX setup.
The PDF is delivered already rendered, so it is not needed to check the
mathematics encoded in the JSON files.

\subsection{Immediate dependency map}
The root comparison pin and baseline recount are recorded under
\texttt{review/}. The named direct calibration interfaces are
\texttt{c222\_iff\_zf\_refutation} and
\texttt{zf\_refutation\_iff\_zfc\_refutation}; the endpoint uses the
selected \texttt{ZFCax\_s} proof predicate. The direct packet retains the exact constructor source,
Horn branch inventory, C222 endpoint, constructible-choice bridge,
Pell/power/factorial and bounded-quantifier interfaces under
\texttt{direct/sources/}. The official Mathlib indexed-Pell completeness
source is fixed independently by its own commit, Git blob and SHA256.
Its source notice and Apache license are preserved.

The radical packet retains the five-row seventeen-integer parent and its
CRT and prime-separation proofs. Its six coefficient payloads are not
interchangeable with a different native tuple. In particular, the
already-shifted $M_F$ must not be shifted again. The ten-variable packet
retains its twelve literal signed integer leaves, their source recipes,
the exact short-guard parent graphs and the parent proof. Included source
inventories distinguish a local dependency from a historical path
mentioned in inherited prose. Historical source references are not
runtime dependencies.

\section{What to pursue next}
\label{sec:next}
The direct quartic is now substantially smaller without changing the
logical endpoint. The remaining sign seeds and the repeated arithmetic
macros are specific targets, but further compression must preserve
simultaneous completeness. Merely orienting a congruence or sharing a
bound because it looks harmless is insufficient.

\begin{enumerate}
\item \textbf{Extend the sign certificate before changing the integer
encoding.} The current $414$ guarded signs, $80$ normalized pivots and $35$
inferred signs establish all $529$ nonquotient semantic signs. A new universally valid implication that removes a seed
saves three unrestricted integer coordinates. It must be checked against
the literal residual and the dependency order; signed quotient cones
must remain outside the natural arithmetic semantics.
\item \textbf{Seek genuinely shared power certificates.} There are still
$40$ positive-power calls. Any proposed shared Pell parameter must be
chosen after all its finite input requirements and must satisfy the
simultaneous index and final-norm conditions. A common large numerical
bound alone does not prove that those Pell witnesses can be shared.
A proof of this joint existence theorem would be more valuable than
another local circuit reassociation.
\item \textbf{Change the short guard or prefix for ten variables.} The
proved $1080888$ lower bound closes the route consisting only of final
polynomial translations. Further degree improvement in that family
requires changing another structural element. The restriction is sharp
and useful, but says nothing about arbitrary ten-variable equations.
\item \textbf{Treat another eliminated guard coordinate as a new theorem.}
The tempting replacement of
$1+p_1^2+p_2^2+p_3(p_3+1)$ by
$1+p_1^2+p_3(p_3+1)$ is not uniformly valid. For $d=1$ and
$\eta\zeta=9$, the positive joint expression is
$9-(t^2-9)^2$. Its only positive values occur at $t=\pm3$ and equal $9$.
But $8$ is not a square plus $z(z+1)$: the only nonnegative consecutive
products at most $8$ are $0,2,6$, leaving $8,6,2$, none a square.
A twelve-variable proposal based on this shortcut therefore needs an
additional canonical-witness tuning argument, not an omitted exception.
This failed shortcut is not the interpolation construction of
Section~\ref{sec:twelve2048}.
\item \textbf{Formalize the smallest complete chain.} A useful next
ProveIt milestone is a checked translation for the direct $1859$-variable
quartic: first the fixed Horn factoring and degree-bootstrap carrier,
then the short-beta and majorant transformations, then the parametric
Pell compiler. The supplied finite checkers identify concrete interfaces
but do not themselves create those universal formal theorems.
\end{enumerate}

\section{Reading and reproducing the archive}
The archive is self-contained at the level claimed here. Its root contains
this report's editable TeX and rendered PDF, a README, the release verifier,
a SHA256 manifest, and a recorded full offline replay. The component
directories contain only the selected equations and their immediate
proof, coefficient and verification dependencies. Earlier complete
incoming archives and redundant development variants are not bundled.

Start with the README, then run \texttt{python verify.py}. This validates
the manifest and executes the full offline check and deterministic
rebuild suite. \texttt{python -O verify.py} provides the corresponding
optimized-driver replay. Outputs made while checking go to an independent
temporary directory; the release itself is left unchanged. To read a
particular equation, use the domain and variant name in the tables above,
not a mixed-domain intermediate template.

The polynomial payloads are the authoritative exact artifacts. The
human-readable proofs explain why they have the required existential
meaning. Neither a successful run nor a small variable count should be
read as a proof that the polynomial has no zero, as global optimality, or
as a substitute for the fixed syntactic definition of consistency.
\end{document}
